\begin{abstract}
Continuous payments let a recipient spend an output before its creating transaction executes publicly. Next lets a recipient verify a quorum-signed output certificate (TXCer) and request a successor after one certification round; public execution follows asynchronously. If the successor executes first, the committee records an obligation against the issuer. Original signers can resubmit a withheld source; otherwise, after the deadline, a public decision debits the issuer's reserve, and a late source repays it. Restricted chameleon-hash commands record the debit in the stored block, preserving authorizations, successor references, and the complete BlockID. For four-member organizations and a four-validator committee with at most one static Byzantine member each, we prove unique consumption, funding of every certificate's outstanding CAL liability, and, for source-closed, conflict-free payment sets, delay-neutral terminal CAL holdings. A certificate does not guarantee public execution, because one global intent admits one payment. With synchronous member commits on one host, 100 hops reach the last certified receipt in a median 4.557\,s versus 64.476\,s with per-hop waiting. In twelve mixed runs (84,192 payments, bounded concurrency), compensation and repayment complete with adaptation paused. Representation lowers local point-read time by 7.44--12.32\% and adds 0.28--2.19\,MiB per repaired block.
\end{abstract}

\begin{IEEEkeywords}
UTXO payments, payment guarantees, Byzantine fault tolerance, authorized historical representation.
\end{IEEEkeywords}